\subsection{赋值除环上的赋范空间}

回忆 (GT, IX, § 3.2)，除环 K 上的绝对值是一个映射 \(\xi \mapsto |\xi|\)（K 到 \(\mathbf{R}_{+}\) 中的映射），使得 \(|\xi| = 0\) 当且仅当 \(\xi = 0\)，并且 \(|\xi \eta| = |\xi| \cdot |\eta|\)，以及 \(|\xi + \eta| \leqslant |\xi| + |\eta|\)；绝对值在 K 上定义了一个距离 \(|\xi - \eta|\)，从而定义了一个与 K 的除环结构相容的 Hausdorff 拓扑。若 \(|\xi| = 1\) 对所有 \(\xi \neq 0\) 成立，则称该绝对值为非真绝对值，它在 K 上定义的拓扑是离散拓扑；反之，若存在 K 中的 \(\alpha \neq 0\) 使 \(|\alpha| \neq 1\)，则存在 K 中的 \(\beta \neq 0\) 使 \(|\beta| < 1\)（只需取 \(\beta = \alpha\) 或 \(\beta = \alpha^{-1}\)），并且序列 \((\beta^n)_{n \geqslant 1}\) 收敛于 0，从而 K 的拓扑不是离散的。

另一方面我们回忆 (GT, IX, § 3.3)，若 E 是非离散赋值除环 K 上的向量空间，则 E 上的范数是一个映射 \(x \mapsto \|x\|\)（E 到 \(\mathbf{R}_{+}\) 中的映射），使得 \(\|x\| = 0\) 当且仅当 \(x = 0\)，并且 \(\| \lambda x \| = |\lambda| \cdot \| x\|\) 对每个标量 \(\lambda \in \mathbf{K}\) 成立，以及 \(\|x + y\| \leqslant \|x\| + \|y\|\)。范数在 E 上定义了一个距离 \(\|x - y\|\)，从而定义了一个与 E 的向量空间结构相容的拓扑 (loc. cit.)。除非另有明确说明，赋范空间总是就其范数所定义的拓扑向量空间结构来考虑。赋范空间属于最重要的拓扑向量空间之列。

我们知道 (GT, IX, § 3.3)，E 上两个不同的范数可以在 E 上定义同一个拓扑；为此必须且只需这两个范数等价 (loc. cit.)。因此赋范空间的结构比拓扑向量空间的结构更为丰富；若 E 和 F 是两个赋范空间，就必须注意区分 E 的赋范空间结构与 F 的赋范空间结构之间的同构概念，以及 E 的拓扑向量空间结构与 F 的拓扑向量空间结构之间的同构概念。

例. --- 设 I 为任意指标集；我们知道 (GT, X, § 3.2)，一个范数 \(\|x\|\) 可以定义在 I 到 K 中的有界映射 \(x = (\xi_{\mathfrak{t}})\) 所成的集合 \(\mathcal{B}(\mathrm{I};\mathrm{K})\)（也记作 \(\mathcal{B}_{\mathrm{K}}(\mathrm{I})\) 或 \(\ell_{\mathrm{K}}^{\infty}(\mathrm{I})\)）上，其定义为 \(\|x\| = \sup_{\mathfrak{t}\in \mathrm{I}}|\xi_{\mathfrak{t}}|\)。当 I 是拓扑空间时，I 到 K 中的有界连续映射所成的集合是空间 \(\mathcal{B}(\mathrm{I};\mathrm{K})\) 的闭子空间 (GT, X, § 3.1, cor. 2)。\(\mathcal{B}(\mathrm{I};\mathrm{K})\) 的另一个子空间是 \(\ell_{\mathrm{K}}^{1}(\mathrm{I})\)，即绝对可和族 \(x = (\xi_{\mathfrak{t}})\) 所成的集合 (GT, X, § 3.6)；在这个子空间上可以定义另一个范数 \(\|x\|_1 = \sum_{\mathfrak{t}\in \mathrm{I}}|\xi_{\mathfrak{t}}|\)，它一般与范数 \(\|x\| = \sup_{\mathfrak{t}\in \mathrm{I}}|\xi_{\mathfrak{t}}|\) 不等价 (I, 第 23 页, 习题 6)；当把 \(\ell_{\mathrm{K}}^{1}(\mathrm{I})\) 视为赋范空间而不指明其范数时，所指的总是范数 \(\|x\|\)。我们用 \(\mathcal{B}(\mathrm{I})\) 与 \(\ell^1 (\mathrm{I})\) 代替 \(\mathcal{B}(\mathrm{I};\mathbf{R})\) 与 \(\ell_{\mathbf{R}}^{1}(\mathrm{I})\)。

